\chapter{Persamaan Diferensial}\label{ch:b2:diffeq}

Tahun ke-1 menyelesaikan persamaan linear yang menerima rumus. Bab ini
memasok apa yang tak dapat diberikan rumus: yakni \emph{teorema Cauchy--Lipschitz}
--- keberadaan dan ketunggalan bagi $y' = f(t, y)$ --- yang dibuktikan lewat
teorema titik tetap Banach, persis seperti dijanjikan pada
\cref{ch:b2:metric}; lalu teori \omterm{def:b2:metric:complete}{lengkap} bagi \emph{sistem
linear} $X' = A(t)X + B(t)$, dengan \omterm{ex:b2:nvs:matrixexp}{eksponensial matriks} dan
\omterm{def:b2:diffeq:wronskian}{wronskian} sebagai mesin hitungnya.

\section{Teorema Cauchy--Lipschitz}

\begin{theorem}[Cauchy--Lipschitz, versi Lipschitz global]\label{thm:b2:diffeq:cauchylipschitz}
Misalkan $I$ sebuah ruas, dan $f \colon I \times \R^n \to \R^n$ \omterm{def:b2:metric:continuity}{kontinu}
serta \emph{\omterm{def:b2:metric:continuity}{Lipschitz} pada peubah keduanya}, secara seragam dalam
peubah pertamanya: yakni $\norm{f(t, y) - f(t, z)} \leq k\,\norm{y - z}$ untuk setiap $t
\in I$. Maka untuk setiap $(t_0, y_0) \in I \times \R^n$, masalah
Cauchy\index{teorema Cauchy--Lipschitz}
\[
y' = f(t, y), \qquad y(t_0) = y_0
\]
punya tepat satu penyelesaian $y \colon I \to \R^n$ berkelas $C^1$.
\end{theorem}

\begin{proof}
\emph{Perumusan ulangnya.} Sebuah $y$ yang \omterm{def:b2:metric:continuity}{kontinu} menyelesaikan masalahnya bila dan hanya bila ia
memenuhi persamaan integral
\[
y(t) = y_0 + \int_{t_0}^{t} f\bigl(s, y(s)\bigr)\,\dd s
=: T(y)(t)
\]
(lewat teorema dasar kalkulus pada kedua arahnya; dan penyelesaian \omterm{def:b2:metric:continuity}{kontinu}
persamaan integralnya otomatis bersifat $C^1$).

\emph{Sebuah kontraksi, setelah dinormakan ulang.} Pada \omterm{def:b2:nvs:banach}{ruang Banach} $E =
C(I, \R^n)$ dengan \omterm{def:b2:nvs:norm}{norma} \emph{terbobot}
\[
N(y) = \sup_{t \in I}\; \eu^{-2k\abs{t - t_0}}\,\norm{y(t)} ,
\]
(yang setara dengan \omterm{def:b2:nvs:norm}{norma} supnya: sebab bobotnya terbatas di atas dan di bawah
pada ruas $I$, jadi $E$ tetap \omterm{def:b2:metric:complete}{lengkap}), taksirlah untuk $y, z \in
E$ dan, katakanlah, $t \geq t_0$:
\[
\norm{T(y)(t) - T(z)(t)}
\leq \int_{t_0}^{t} k\,\norm{y(s) - z(s)}\,\dd s
\leq k\,N(y - z)\int_{t_0}^{t} \eu^{2k(s - t_0)}\dd s
\leq \frac{N(y-z)}{2}\,\eu^{2k(t - t_0)} .
\]
Mengalikannya dengan $\eu^{-2k(t - t_0)}$ lalu mengambil supnya (dengan kasus $t
< t_0$ yang simetris): $N\bigl(T(y) - T(z)\bigr) \leq \frac12 N(y -
z)$: jadi $T$ merupakan kontraksi berfaktor $\frac12$ pada $(E, N)$ yang \omterm{def:b2:metric:complete}{lengkap}. Lalu
teorema titik tetap Banach (\cref{thm:b2:metric:banach}) menghasilkan
titik tetap yang tunggal: yakni penyelesaian tunggalnya.
\end{proof}

\begin{remark}
Untuk $f$ yang sekadar $C^1$ (yakni \omterm{def:b2:metric:continuity}{Lipschitz} secara lokal), teoremanya berlaku
secara \emph{lokal}, dengan penyelesaian maksimal pada interval \omterm{def:b2:metric:topology}{terbuka} maksimalnya;
dan penyelesaiannya dapat meledak dalam waktu berhingga ($y' = y^2$, $y(0) = 1$: $y(t)
= \frac{1}{1-t}$, yang lenyap di $t = 1$). Jadi hipotesis \omterm{def:b2:metric:continuity}{Lipschitz} globalnya
yang membeli seluruh ruasnya. Dua akibat yang layak dipahat:
kurva penyelesaian sebuah persamaan \omterm{def:b2:diffcalc:differential}{diferensial} dengan medan \omterm{def:b2:metric:continuity}{Lipschitz} \emph{tak pernah bersilangan};
dan fungsi nol adalah satu-satunya penyelesaian yang lenyap di mana pun bagi
persamaan homogen yang linear.
\end{remark}

\begin{example}[Ketunggalan itu sebuah teorema: sebuah medan yang bocor]\label{ex:b2:diffeq:leaky}
Tinjaulah $y' = 2\sqrt{\abs y}$ dengan $y(0) = 0$. Fungsi nol
menyelesaikannya; demikian pula
\[
y(t) = \begin{cases} 0 & t \leq 0,\\ t^2 & t \geq 0,
\end{cases}
\]
yang bersifat $C^1$ (sebab kedua kepingnya berturunan $0$ di titik
lemnya) dan memenuhi $y'(t) = 2t = 2\sqrt{t^2}$ untuk $t > 0$ ---
bahkan, menunda lepas landasnya memberi penyelesaian bagi \emph{setiap}
waktu pelepasan $c \geq 0$: jadi tak berhingga banyak penyelesaian lewat
data awal yang sama. Dan tak ada pertentangan dengan
\cref{thm:b2:diffeq:cauchylipschitz}: sebab di dekat $y = 0$,
\[
\frac{\abs{2\sqrt y - 2\sqrt z}}{\abs{y - z}}
= \frac{2}{\sqrt y + \sqrt z} \longrightarrow +\infty ,
\]
medannya tidak \omterm{def:b2:metric:continuity}{Lipschitz} dalam $y$, jadi teoremanya membisu.
Pelajaran penutupnya: pembacaan fisisnya adalah sebuah ember yang mengering
di bawah gravitasi lalu dijalankan mundur --- sebab dari keadaan kosongnya, kita
tak dapat mengetahui kapan ia mulai terisi; jadi determinisme persamaan \omterm{def:b2:diffcalc:differential}{diferensial}
persis merupakan syarat \omterm{def:b2:metric:continuity}{Lipschitznya}, bukan hukum alam.
\end{example}

\section{Sistem linear}

\begin{theorem}[Struktur sistem linear]\label{thm:b2:diffeq:linear}
Misalkan $A \colon I \to \mathcal{M}_n(\R)$ dan $B \colon I \to \R^n$
\omterm{def:b2:metric:continuity}{kontinu} pada interval $I$. Untuk setiap $(t_0, X_0)$ masalah
\[
X' = A(t)X + B(t), \qquad X(t_0) = X_0
\]
punya tepat satu penyelesaian pada \emph{seluruh} $I$. Adapun penyelesaian
sistem homogennya (dengan $B = 0$) membentuk ruang vektor $\mathcal{S}_H$ yang
berdimensi tepat $n$, dan penilaian $X \mapsto X(t_0)$ merupakan
isomorfisma $\mathcal{S}_H \to \R^n$; jadi penyelesaian umum $=$
penyelesaian khusus $+$ penyelesaian homogen.
\end{theorem}

\begin{proof}
Pada setiap ruas $J \subseteq I$ yang memuat $t_0$: pemetaan $f(t, X) =
A(t)X + B(t)$ bersifat \omterm{def:b2:metric:continuity}{kontinu}, dan \omterm{def:b2:metric:continuity}{Lipschitz} dalam $X$ dengan konstanta $k
= \sup_J \vertiii{A(t)}$ (yang berhingga sebab \omterm{def:b2:metric:continuity}{kontinu} pada sebuah ruas):
jadi \cref{thm:b2:diffeq:cauchylipschitz} berlaku pada $J$; lalu dengan membiarkan $J$
menghabiskan $I$, ketunggalannya melem penyelesaiannya menjadi satu pada $I$.
Adapun kelinearan himpunan penyelesaiannya dan pemetaan penilaiannya sudah jelas;
sedangkan penilaiannya bijektif berkat keberadaannya (jadi surjektif) dan ketunggalannya
(jadi injektif): sehingga $\dim \mathcal{S}_H = n$. Adapun struktur afinnya adalah hujah
Tahun ke-1 kata demi kata.
\end{proof}

\begin{example}[Isomorfisma penilaian, secara konkret]\label{ex:b2:diffeq:evaliso}
Untuk $y'' + y = 0$, yang dipandang sebagai sistem $X' =
\begin{pmatrix} 0 & 1\\ -1 & 0\end{pmatrix}X$ dengan $X = (y,
y')$: teoremanya mengatakan bahwa ruang penyelesaiannya sebuah bidang, dan bahwa
$X \mapsto X(0) = (y(0), y'(0))$ merupakan isomorfisma ke
$\R^2$. Adapun penyelesaian $\cos$ dan $\sin$ bernilai $(1, 0)$
dan $(0, 1)$ --- yakni basis kanonik $\R^2$ --- jadi keduanya membentuk
basis ruang penyelesaiannya, dan \emph{setiap} penyelesaiannya berbunyi
\[
y(t) = y(0)\cos t + y'(0)\sin t ,
\]
dengan koefisiennya dibaca langsung dari data awalnya, tanpa
sistem linear yang harus diselesaikan. Pelajaran penutupnya: memilih
\omterm{def:b2:diffeq:wronskian}{sistem fundamental} yang nilai awalnya berupa basis
kanoniknya (di sini $\cos, \sin$) persis merupakan memilih kolom
$\eu^{tA}$; jadi isomorfisma penilaiannya adalah sebab syarat
awalnya memparameterkan lintasannya --- yakni isi geometris
``dinamika deterministik'' bagi persamaan linear.
\end{example}

\begin{definition}[Wronskian]\label{def:b2:diffeq:wronskian}
Untuk penyelesaian $X_1, \dots, X_n$ sistem homogennya,
\emph{wronskian}\index{wronskian} adalah $W(t) = \det\bigl(X_1(t),
\dots, X_n(t)\bigr)$. Berkat isomorfisma di atas, entah $W$ lenyap
secara identik (jadi keluarganya bergantungan) atau tak pernah lenyap (jadi sebuah \emph{sistem
fundamental}\index{sistem fundamental}); dan secara kuantitatif, $W' = \operatorname{tr}\bigl(A(t)\bigr)
W$, sehingga
\[
W(t) = W(t_0)\,\exp\Bigl(\int_{t_0}^{t}
\operatorname{tr} A(s)\,\dd s\Bigr)
\quad \text{(rumus Liouville)}.
\]
\end{definition}

\begin{example}[Liouville diperiksa pada sebuah persamaan Euler]\label{ex:b2:diffeq:eulerwronskian}
Pada $\intoo{0}{\infty}$, persamaan $t^2y'' + ty' - y = 0$ punya
penyelesaian $y_1(t) = t$ dan $y_2(t) = \frac1t$ (sulihkan).
Adapun \omterm{def:b2:diffeq:wronskian}{wronskiannya}:
\[
W(t) = \det\begin{pmatrix} t & \tfrac1t\\[2pt]
1 & -\tfrac{1}{t^2}\end{pmatrix}
= -\frac1t - \frac1t = -\frac2t ,
\]
yang tak pernah nol: jadi sebuah \omterm{def:b2:diffeq:wronskian}{sistem fundamental}. Kini periksalah Liouville: dalam
bentuk ternormalkan $y'' + \frac1t\,y' - \frac{1}{t^2}\,y = 0$,
matriks pendampingnya $A(t) = \begin{pmatrix} 0 & 1\\ \frac{1}{t^2}
& -\frac1t\end{pmatrix}$ bertrace $-\frac1t$, jadi
\[
W(t) = W(1)\exp\Bigl(-\int_1^t\frac{\dd s}{s}\Bigr)
= -2\,\eu^{-\ln t} = -\frac2t . \checkmark
\]
Pelajaran penutupnya: Liouville meramalkan \emph{bentuk}
\omterm{def:b2:diffeq:wronskian}{wronskiannya} sebelum satu penyelesaian pun diketahui --- di sini, bahwa $W$ mestilah
$\frac{c}{t}$; dan inilah yang menggerakkan metode
penurunan orde
(\cref{prop:b2:diffeq:secondorder}), tempat mengetahui $y_1$ dan
bentuk \omterm{def:b2:diffeq:wronskian}{wronskiannya} menentukan $y_2$ lewat satu kuadratur.
\end{example}

\begin{proof}[Bukti rumus Liouville]
Berlaku $W(t) = \det M(t)$ dengan $M' = AM$. Menurunkan \omterm{def:b2:linalg:det}{determinannya}
sebagai fungsi multilinear atas kolomnya memberi
\[
W'(t) = \sum_j \det(X_1, \dots, X_j', \dots, X_n)
= \sum_j \det(X_1, \dots, AX_j, \dots, X_n) .
\]
Kini pemetaan $(C_1, \dots, C_n) \mapsto \sum_j \det(C_1, \dots,
AC_j, \dots, C_n)$ bersifat $n$-linear dan berselang-seling (sebab dengan dua kolom yang sama
$C_i = C_k$, suku dengan $j \notin \{i, k\}$ lenyap seketika,
sedangkan suku $j = i$ dan $j = k$ saling meniadakan berpasangan setelah satu
pertukaran kolom): jadi menurut teorema ketunggalannya
(\cref{thm:b2:linalg:detspace}) ia sama dengan $c \cdot \det$, dengan $c$ dibaca
pada kolom kanoniknya: $c = \sum_j \det(e_1, \dots, Ae_j, \dots,
e_n) = \sum_j a_{jj} = \operatorname{tr} A$. Karena itu $W' =
\operatorname{tr}\bigl(A(t)\bigr)W$: yakni persamaan \omterm{def:b2:diffcalc:differential}{diferensial} linear skalar, yang diselesaikan
lewat rumus Tahun ke-1.
\end{proof}

\section{Koefisien tetap: eksponensial matriks}

\begin{theorem}\label{thm:b2:diffeq:matrixexp}
Untuk $A \in \mathcal{M}_n(\R)$ (atau $\C$), eksponensialnya $\eu^{tA}
= \sum_k \frac{(tA)^k}{k!}$ (\cref{ex:b2:nvs:matrixexp}) memenuhi:
bahwa $t \mapsto \eu^{tA}$ bersifat $C^1$ (bahkan $C^\infty$)
dengan\index{eksponensial matriks}
\[
\frac{\dd}{\dd t}\,\eu^{tA} = A\,\eu^{tA} = \eu^{tA}A ,
\qquad
\eu^{(s+t)A} = \eu^{sA}\,\eu^{tA},
\qquad
(\eu^{A})^{-1} = \eu^{-A} ;
\]
dan $\eu^{A + B} = \eu^A\eu^B$ \emph{apabila} $AB = BA$. Adapun masalah
Cauchy $X' = AX$ dengan $X(0) = X_0$ punya penyelesaian tunggal $X(t) =
\eu^{tA}X_0$; dan dengan sumber, rumus \emph{variasi konstanta}
berlaku:
\[
X(t) = \eu^{(t - t_0)A}X_0 + \int_{t_0}^{t} \eu^{(t-s)A}B(s)\,\dd
s .
\]
\end{theorem}

\begin{proof}
\emph{Keterdiferensialannya:} deret $\sum \frac{t^kA^k}{k!}$ dan
deret turunan sukunya $\sum \frac{t^{k-1}A^k}{(k-1)!} = A\sum
\frac{(tA)^{k-1}}{(k-1)!}$ konvergen normal pada setiap ruas
(dengan \omterm{def:b2:nvs:norm}{norma} $\leq \frac{(\abs t\,\vertiii A)^k}{k!}$): jadi turunkan
suku demi suku (\cref{thm:b2:funcseq:seriestransfer}, versi bernilai vektor).
Adapun kedua urutannya $A\eu^{tA}$ dan $\eu^{tA}A$ sepakat karena setiap
jumlah parsialnya komut dengan $A$.

\emph{Hukum grupnya:} untuk $A, B$ yang komut, hasil kali Cauchy kedua
deret eksponensialnya ditata ulang lewat teorema binomial persis
seperti pada \cref{ex:b2:series:exp} (sebab kekonvergenan \omterm{def:b2:series:def}{mutlak} pada aljabar
Banachnya membenarkannya): jadi $\eu^{A+B} = \eu^A\eu^B$; lalu dengan $B = sA$ ini
memberi hukum grup satu parameternya, dan dengan $B = -A$ balikannya.

\emph{Masalah Cauchynya:} $X(t) = \eu^{tA}X_0$ menyelesaikannya (turunkan saja);
dan ketunggalannya berkat \cref{thm:b2:diffeq:linear}. Adapun variasi konstantanya:
tetapkan $Y(t) = \eu^{-tA}X(t)$; lalu setelah diturunkan, $Y' = \eu^{-tA}(X' -
AX) = \eu^{-tA}B(t)$; jadi integralkan dari $t_0$ ke $t$ lalu kalikan
kembali dengan $\eu^{tA}$.
\end{proof}

\begin{method}[Menghitung $\eu^{tA}$]\label{met:b2:diffeq:computeexp}
Reduksikan $A$ (\cref{ch:b2:reduction}): bila $A = PDP^{-1}$ berbentuk diagonal,
maka $\eu^{tA} = P\,\eu^{tD}P^{-1}$ dengan $\eu^{tD}$ diagonal berisi
$\eu^{t\lambda_i}$; sedangkan secara umum pakailah Dunford $A = D + N$
(yang komut): $\eu^{tA} = \eu^{tD}\,\eu^{tN}$ dengan $\eu^{tN}$ sebuah
\emph{polinomial} dalam $t$ (sebab kenilpotenannya memotong deretnya). Adapun \omterm{def:b2:reduction:eigen}{nilai
eigen} kompleksnya berpasangan menjadi blok perputaran-kali-eksponensial
(\cref{exo:b2:diffeq:5}).
\end{method}

\begin{remark}[Jebakan yang sering muncul]
\emph{(i) $\eu^{A+B} \neq \eu^A\eu^B$ tanpa kekomutan:}
ambillah $A = \begin{pmatrix} 0 & 1\\ 0 & 0\end{pmatrix}$ dan $B =
\begin{pmatrix} 0 & 0\\ 1 & 0\end{pmatrix}$. Maka $\eu^A = I +
A$ dan $\eu^B = I + B$ (berkat kenilpotenannya), jadi
\[
\eu^A\eu^B = \begin{pmatrix} 2 & 1\\ 1 & 1\end{pmatrix},
\qquad\text{sedangkan}\qquad
\eu^{A+B} = \cosh(1)\,I + \sinh(1)\,(A + B)
= \begin{pmatrix} \cosh 1 & \sinh 1\\ \sinh 1 & \cosh 1
\end{pmatrix},
\]
dengan memakai $(A+B)^2 = I$; dan $\cosh 1 \approx 1.54 \neq 2$. Jadi
hukum grup pada \cref{thm:b2:diffeq:matrixexp} mengusung hipotesis
yang sejati. \emph{(ii) Gerak hati nonlinear di lahan linear:}
penyelesaian sistem yang \emph{linear} dengan koefisien
\omterm{def:b2:metric:continuity}{kontinu} hidup pada seluruh intervalnya
(\cref{thm:b2:diffeq:linear}) --- jadi bila sebuah calon penyelesaian
meledak di dalam $I$, maka persamaannya tak linear atau
perhitungannya salah; sebaliknya, untuk persamaan nonlinear
janganlah pernah menjanjikan kegloballan tanpa hujah ($y' = y^2$).
\emph{(iii) Membagi dengan yang tak diketahui:} memisahkan peubah pada
$y' = y(1-y)$ diam-diam membuang penyelesaian tetapnya $0$ dan
$1$ --- yakni persis yang menata garis fasanya
(\cref{exo:b2:diffeq:3}); jadi daftarkanlah penyelesaian tetapnya lebih dulu.
\emph{(iv) Data awalnya menetapkan vektor, bukan skalar:} sebuah persamaan
skalar berorde $n$ menuntut $n$ syarat (yakni $y, y', \dots$ di
$t_0$); jadi mencocokkan $y(t_0)$ saja menyisakan keluarga
berparameter $(n-1)$, yakni sumber klasik bagi konstanta yang ``hilang''.
\end{remark}

\begin{example}[Sebuah eksponensial $3\times3$ lewat Dunford]\label{ex:b2:diffeq:dunford3}
Selesaikan $X' = AX$ untuk $A = \begin{pmatrix} 2 & 1 & 0\\ 0 & 2 &
0\\ 0 & 0 & 3\end{pmatrix}$. Dunford blok demi blok: $A = D + N$
dengan $D = \operatorname{diag}(2, 2, 3)$ dan $N = E_{12}$, yang
komut (sebab $N$ hidup di dalam blok bernilai eigen $2$), dan $N^2 =
0$:
\[
\eu^{tA} = \eu^{tD}\,\eu^{tN}
= \begin{pmatrix}
\eu^{2t} & t\,\eu^{2t} & 0\\
0 & \eu^{2t} & 0\\
0 & 0 & \eu^{3t}
\end{pmatrix} .
\]
Penyelesaian umumnya terbaca kolom demi kolom: $X(t) =
\bigl(\eu^{2t}(x_0 + ty_0),\ \eu^{2t}y_0,\
\eu^{3t}z_0\bigr)$. Pemeriksaan kewarasannya: di $t = 0$ matriksnya adalah
$I$; \omterm{def:b2:linalg:det}{determinannya} $\eu^{7t} =
\eu^{t\operatorname{tr}A}$, seperti dituntut Liouville; dan
faktor $t$-nya muncul persis di tempat \omterm{def:b2:reduction:eigen}{nilai eigen} $2$-nya
cacat. Pelajaran penutupnya: polinomial kali eksponensial itu
bukan tebakan yang harus dihafal --- sebab ia adalah deret terpotong
$\eu^{tN}$, dan derajatnya dibatasi indeks kenilpotenannya,
tak pernah lebih.
\end{example}

\begin{method}[Menyelesaikan $X' = AX + B(t)$, dari awal sampai akhir]\label{met:b2:diffeq:solvelinear}
\begin{enumerate}
  \item \omterm{def:b2:reduction:eigen}{Spektrum} $A$; lalu $\eu^{tA}$ lewat
        \cref{met:b2:diffeq:computeexp} (yakni diagonalkan; atau
        lewat Dunford seperti pada \cref{ex:b2:diffeq:dunford3}; atau lewat
        siasat polinomial seperti $A^2 = -I$).
  \item Sebuah penyelesaian khusus: variasi konstantanya
        $\int_{t_0}^t\eu^{(t-s)A}B(s)\dd s$ selalu berhasil;
        sedangkan untuk $B$ yang eksponensial-polinomial, ansatz berbentuk
        sama (dengan derajat yang dinaikkan pada \omterm{pb:b2:diffeq:1}{resonansi},
        \cref{exo:b2:diffeq:10}) lebih cepat.
  \item Penyelesaian umum $= \eu^{(t-t_0)A}X_0 +$ penyelesaian khususnya;
        lalu cocokkan data awalnya \emph{paling akhir}, pada
        rumus yang \omterm{def:b2:metric:complete}{lengkap}.
  \item Pemeriksaan kewarasannya: $X(t_0)$ benar; pertumbuhan bagian
        homogennya cocok dengan bagian real \omterm{def:b2:reduction:eigen}{nilai eigennya}
        (\cref{exo:b2:diffeq:8}); dan $\det$ sebuah matriks
        fundamentalnya menaati Liouville.
\end{enumerate}
\end{method}

\begin{example}[Sebuah potret fasa]\label{ex:b2:diffeq:phaseportrait}
$X' = AX$ dengan $A = \begin{pmatrix} 0 & 1\\ -1 & 0\end{pmatrix}$:
karena $A^2 = -I$, deretnya terpecah menjadi
\[
\eu^{tA} = (\cos t)\,I + (\sin t)\,A
= \begin{pmatrix} \cos t & \sin t\\ -\sin t & \cos t
\end{pmatrix} :
\]
jadi lintasannya berupa lingkaran yang ditempuh searah jarum jam --- yakni osilator harmonik
$x'' + x = 0$ dalam busana orde pertama. \omterm{def:b2:reduction:eigen}{Nilai eigennya} $\pm\iu$ pada
sumbu imajinernya: jadi sebuah \emph{pusat}. Lebih umum, bagian real
\omterm{def:b2:reduction:eigen}{nilai eigen} $A$ memutuskan pertumbuhan atau peluruhan $\norm{X(t)}$
(\cref{exo:b2:diffeq:8}).
\end{example}

\begin{omfigure}
\begin{tikzpicture}
\begin{axis}[omaxis, width=6.4cm, height=6.4cm,
    xmin=-2.3, xmax=2.3, ymin=-2.3, ymax=2.3,
    xtick={-1,1}, ytick={-1,1}]
  \addplot[omDef, thick, domain=0:360, samples=90, smooth]
    ({cos(x)}, {-sin(x)});
  \addplot[omDef!60, thick, domain=0:360, samples=90, smooth]
    ({1.7*cos(x)}, {-1.7*sin(x)});
  \draw[-Stealth, omThm, thick] (axis cs:1,0) -- (axis cs:0.995,-0.12);
  \draw[-Stealth, omThm, thick] (axis cs:1.7,0) -- (axis cs:1.695,-0.12);
\end{axis}
\end{tikzpicture}
\hspace{6mm}
\begin{tikzpicture}
\begin{axis}[omaxis, width=6.4cm, height=6.4cm,
    xmin=-2.3, xmax=2.3, ymin=-2.3, ymax=2.3,
    xtick={-1,1}, ytick={-1,1}]
  % A = diag(-1,-2)-type stable node: trajectories y = c x^2
  \addplot[omMeth, thick, domain=-1.5:1.5, samples=60] {x^2};
  \addplot[omMeth, thick, domain=-1.5:1.5, samples=60] {-x^2};
  \addplot[omMeth, thick, domain=-2.2:2.2, samples=2] {0};
  \addplot[omMeth!60, thick, domain=-1.2:1.2, samples=60] {2*x^2};
  \draw[-Stealth, omThm, thick] (axis cs:1.2,1.44) --
    (axis cs:1.05,1.1);
  \draw[-Stealth, omThm, thick] (axis cs:-1.2,-1.44) --
    (axis cs:-1.05,-1.1);
\end{axis}
\end{tikzpicture}

{\small Dua potret fasa yang linear. Kiri: sebuah \emph{pusat}
(dengan \omterm{def:b2:reduction:eigen}{nilai eigen} $\pm\iu$) --- yakni orbit lingkaran tertutup milik osilator
harmonik. Kanan: sebuah \emph{simpul stabil} (dengan \omterm{def:b2:reduction:eigen}{nilai eigen} $-1, -2$) ---
tempat semua lintasannya jatuh ke titik asalnya secara menyinggung arah eigen
yang lambat.}
\end{omfigure}

\begin{omfigure}
\begin{tikzpicture}
\begin{axis}[omaxis, width=9.6cm, height=6.4cm,
    xmin=-3.2, xmax=3.2, ymin=-1.6, ymax=2.6,
    xtick={0}, ytick={0},
    xlabel={$\tau = \operatorname{tr}A$},
    ylabel={$\delta = \det A$}]
  \addplot[omThm, very thick, domain=-3:3, samples=100]
    {x^2/4};
  \node[font=\small, omThm] at (axis cs:2.45,2.1)
    {$\delta = \tau^2/4$};
  \node[font=\small] at (axis cs:0,-0.9) {pelana};
  \node[font=\small] at (axis cs:-2.45,0.7)
    {simpul stabil};
  \node[font=\small] at (axis cs:2.45,0.7)
    {simpul takstabil};
  \node[font=\small] at (axis cs:-1.05,1.9)
    {spiral stabil};
  \node[font=\small] at (axis cs:1.05,1.9)
    {spiral takstabil};
  \node[font=\small, omDef] at (axis cs:0,2.35) {pusat};
  \draw[omDef, very thick] (axis cs:0,0) -- (axis cs:0,2.55);
\end{axis}
\end{tikzpicture}

{\small Bidang trace--determinan bagi $X' = AX$ pada dimensi
$2$: di bawah sumbu mendatarnya, pelana; di antara sumbunya dan
parabola $\delta = \tau^2/4$, simpul; di dalam parabolanya,
spiral; dan pada sumbu-$\delta$ yang positif, pusat. Adapun soal akhir
pekan membuktikan penggolongan ini lalu menyusuri satu garis tegak
darinya --- yakni osilator teredam --- sampai ke \omterm{pb:b2:diffeq:1}{resonansi}.}
\end{omfigure}

\begin{remark}[Di mana ini dipakai]
Sistem linear adalah model lokal bagi segala yang nonlinear:
sebab di dekat kesetimbangannya, medan vektor yang mulus berkelakuan (pada
kasus hiperboliknya) seperti pelinearannya, yang potretnya digolongkan
bidang trace--determinan. Adapun soal akhir pekan mengerjakan
cerita osilatornya selengkapnya --- peredaman, pemaksaan,
\omterm{pb:b2:diffeq:1}{resonansi}, dan teorema pembandingan Sturm bagi koefisien
yang berubah --- yakni matematika di balik peredam kejut, rangkaian
arus bolak-balik, dan senjang spektral sekaligus. Adapun jilid Tahun ke-3 kembali
dengan teori kualitatifnya (aliran, kestabilan, integral
pertama) pada manifold.
\end{remark}

\section{Orde dua dengan koefisien yang berubah}

\begin{proposition}\label{prop:b2:diffeq:secondorder}
Persamaan $y'' + a(t)y' + b(t)y = c(t)$ (dengan $a, b, c$ \omterm{def:b2:metric:continuity}{kontinu} pada
$I$) adalah sistem $X' = A(t)X + B(t)$ untuk $X = (y, y')$: jadi penyelesaiannya
ada dan tunggal pada seluruh $I$ bagi sembarang data awal $(y(t_0),
y'(t_0))$; dan penyelesaian homogennya membentuk sebuah bidang. Bila satu
penyelesaian homogen $y_1$ yang tak lenyap diketahui, maka penyelesaian bebas keduanya
ditemukan lewat \emph{penurunan orde}: sebab menetapkan $y = y_1 z$ mengubah
persamaan homogennya menjadi persamaan orde pertama bagi $z'$, yang diselesaikan
lewat kuadratur.
\end{proposition}

\begin{proof}
Bentuk sistemnya beserta \cref{thm:b2:diffeq:linear} memberi segala yang
struktural. Adapun penurunannya: setelah menyulihkan $y = y_1z$,
\[
y_1 z'' + (2y_1' + a y_1)z' + \underbrace{(y_1'' + ay_1' +
by_1)}_{=\,0}\,z = 0 :
\]
yakni persamaan linear orde pertama dalam $u = z'$, yang terselesaikan lewat rumus
Tahun ke-1; lalu mengintegralkan $u$ memberi $z$, sehingga $y_2 = y_1 z$,
yang bebas dari $y_1$ setiap kali $z$ tak tetap.
\end{proof}

\begin{example}\label{ex:b2:diffeq:lowering}
$t^2y'' - 2y = 0$ pada $\intoo{0}{\infty}$: fungsi $y_1 = t^2$ merupakan sebuah
penyelesaian. Sulihkan $y = t^2z$: dari $y' = t^2z' + 2tz$ dan $y''
= t^2z'' + 4tz' + 2z$,
\[
t^2y'' - 2y = t^4 z'' + 4t^3z' = 0,
\qquad\text{yakni}\qquad \frac{z''}{z'} = -\frac4t :
\]
jadi $z' = t^{-4}$ (hingga sebuah konstanta), $z = -\frac{1}{3t^3}$, dan $y_2
= t^2z = -\frac{1}{3t}$. Penyelesaian umumnya: $y = \alpha t^2 +
\frac{\beta}{t}$.
\end{example}

\section{Latihan}

\begin{exercise}[$\star$]\label{exo:b2:diffeq:1}
Selesaikan $X' = AX$ dengan $X(0) = (1, 0)^{\mathsf T}$, untuk $A =
\begin{pmatrix} 1 & 1\\ 0 & 2\end{pmatrix}$ (lewat pendiagonalan) dan $A =
\begin{pmatrix} 2 & 1\\ 0 & 2 \end{pmatrix}$ (lewat Dunford).
\end{exercise}

\begin{exercise}[$\star$]\label{exo:b2:diffeq:2}
Masalah Cauchy mana yang punya penyelesaian global tunggal pada $\R$ menurut
\cref{thm:b2:diffeq:cauchylipschitz}? $y' = \sin(ty)$; $\;y' =
y^2$; $\;y' = \abs y$. Untuk yang terakhir, selesaikanlah secara gamblang dengan $y(0) =
0$ dan $y(0) = 1$.
\end{exercise}

\begin{exercise}[$\star$]\label{exo:b2:diffeq:3}
Buktikan bahwa dua penyelesaian maksimal yang berbeda bagi $y' = f(t,y)$ (dengan $f$
\omterm{def:b2:metric:continuity}{Lipschitz} dalam $y$) tak pernah bernilai sama pada waktu yang sama, lalu
turunkan bahwa penyelesaian $y' = y(1 - y)$ yang bermula di
$\intoo{0}{1}$ tinggal di $\intoo{0}{1}$ selamanya.
\end{exercise}

\begin{exercise}[$\star\star$]\label{exo:b2:diffeq:4}
Hitunglah $\eu^{tA}$ untuk $A = \begin{pmatrix} 3 & 1\\ -1 &
1\end{pmatrix}$ \emph{(lewat Dunford: $(A - 2I)^2 = 0$)}, lalu selesaikan $X' =
AX + \begin{pmatrix} \eu^{2t}\\ 0\end{pmatrix}$ dengan $X(0) = 0$, lewat
variasi konstanta.
\end{exercise}

\begin{exercise}[$\star\star$]\label{exo:b2:diffeq:5}
Untuk $A = \begin{pmatrix} \alpha & -\beta\\ \beta &
\alpha\end{pmatrix}$, buktikan $\eu^{tA} =
\eu^{\alpha t}\begin{pmatrix} \cos\beta t & -\sin\beta t\\
\sin\beta t & \cos\beta t\end{pmatrix}$ --- yakni lintasan spiral ---
lewat dua cara: lewat deretnya (tulislah $A = \alpha I + \beta J$ dengan $J^2 =
-I$), dan lewat pengenalan kompleks $z' = (\alpha +
\iu\beta)z$.
\end{exercise}

\begin{exercise}[$\star\star$]\label{exo:b2:diffeq:6}
(Lema Gronwall) Misalkan $u$ \omterm{def:b2:metric:continuity}{kontinu} dan taknegatif dengan $u(t)
\leq C + k\int_{t_0}^{t} u(s)\,\dd s$ pada $\intco{t_0}{T}$. Buktikan
$u(t) \leq C\,\eu^{k(t - t_0)}$ \emph{(turunkan $v(t) =
\eu^{-kt}\int_{t_0}^t u$)}. Lalu turunkan sekali lagi ketunggalan pada
Cauchy--Lipschitz beserta kebergantungan \omterm{def:b2:metric:continuity}{kontinu} $\norm{y(t) -
z(t)} \leq \norm{y_0 - z_0}\,\eu^{k\abs{t - t_0}}$ bagi dua
penyelesaian dengan data awal yang berbeda.
\end{exercise}

\begin{exercise}[$\star\star$]\label{exo:b2:diffeq:7}
Dengan mengetahui bahwa $y_1(t) = \frac{\sin t}{t}$ menyelesaikan $ty'' + 2y' + ty =
0$ pada $\intoo{0}{\pi}$, carilah penyelesaian bebas keduanya lewat
penurunan orde, lalu berikanlah penyelesaian umumnya.
\end{exercise}

\begin{exercise}[$\star\star\star$]\label{exo:b2:diffeq:8}
Misalkan $A \in \mathcal{M}_n(\C)$ dengan semua \omterm{def:b2:reduction:eigen}{nilai eigennya} berbagian real
negatif (secara tegas). Buktikan bahwa setiap penyelesaian $X' = AX$ menuju
$0$ ketika $t \to +\infty$, dengan laju eksponensial: yakni $\norm{X(t)}
\leq C\,\eu^{-\alpha t}$ untuk suatu $\alpha > 0$.
\emph{(Trigonalkan; lalu tanganilah sistem segitiganya dari baris terakhir
ke atas, atau pakailah Dunford: $\eu^{tA} = \eu^{tD}\eu^{tN}$ dengan
$\norm{\eu^{tD}} \leq \eu^{-\alpha' t}$ dan $\eu^{tN}$ polinomial
dalam $t$.)}
\end{exercise}

\begin{exercise}[$\star\star\star$]\label{exo:b2:diffeq:9}
(Tak ada pelarian dalam waktu berhingga bagi pertumbuhan linear) Andaikan $f$
\omterm{def:b2:metric:continuity}{kontinu} dengan $\norm{f(t, y)} \leq a\norm y + b$ pada
$\intco{0}{\infty} \times \R^n$, dan \omterm{def:b2:metric:continuity}{Lipschitz} secara lokal dalam $y$. Dengan memakai
Gronwall (\cref{exo:b2:diffeq:6}) pada bentuk integralnya, buktikanlah bahwa
penyelesaian maksimalnya bersifat global (yakni terdefinisi pada seluruh
$\intco{0}{\infty}$).
\end{exercise}

\begin{exercise}[$\star$]\label{exo:b2:diffeq:10}
Selesaikan $y'' - 3y' + 2y = \eu^{t}$: yakni penyelesaian homogennya, lalu sebuah
penyelesaian khusus berbentuk $\alpha t\,\eu^{t}$ \emph{(mengapa
tebakan naif $\alpha\eu^t$ gagal?)}; lalu penyelesaian umumnya dan
penyelesaian dengan $y(0) = y'(0) = 0$.
\end{exercise}

\begin{exercise}[$\star\star$]\label{exo:b2:diffeq:11}
Hitunglah $\eu^{tA}$ untuk blok Jordan
\[
A = \begin{pmatrix} \lambda & 1 & 0\\ 0 & \lambda & 1\\
0 & 0 & \lambda\end{pmatrix},
\]
lalu paparkanlah semua penyelesaian $X' = AX$: yakni eksponensial kali
vektor polinomial, dengan derajat sampai $2$. Dari mana derajat
polinomialnya datang?
\end{exercise}

\begin{exercise}[$\star\star\star$]\label{exo:b2:diffeq:12}
(Pemaksaan berkala, tanggapan berkala) Misalkan $A \in
\mathcal{M}_n(\R)$ dan $B \colon \R \to \R^n$ \omterm{def:b2:metric:continuity}{kontinu} serta
berkala-$T$.
\begin{enumerate}
  \item Tunjukkan bahwa sebuah penyelesaian $X' = AX + B(t)$ bersifat
        berkala-$T$ bila dan hanya bila $X(T) = X(0)$
        \emph{(bandingkan $X(\cdot + T)$ dan $X$)}.
  \item Tunjukkan bahwa \omterm{def:b2:reduction:eigen}{nilai eigen} $\eu^{TA}$ adalah
        $\eu^{T\lambda}$ dengan $\lambda \in \operatorname{Sp}A$
        \emph{(trigonalkan atas $\C$)}. Lalu turunkan: bahwa bila tak ada
        \omterm{def:b2:reduction:eigen}{nilai eigen} $A$ yang terletak di $\frac{2\iu\pi}{T}\Z$, maka
        $I - \eu^{TA}$ bersifat terbalikkan.
  \item Di bawah hipotesis itu, buktikan bahwa sistemnya punya
        tepat satu penyelesaian berkala-$T$, dengan
        \[
        X(0) = \bigl(I - \eu^{TA}\bigr)^{-1}
        \int_0^{T}\eu^{(T-s)A}B(s)\,\dd s .
        \]
        Adapun kasus yang dikecualikan itu bersesuaian dengan apa, bagi
        osilator harmoniknya? (Soal akhir pekan menjawabnya:
        yakni \omterm{pb:b2:diffeq:1}{resonansi}.)
\end{enumerate}
\end{exercise}

\section{Soal: Osilasi, resonansi, dan teorema pembandingan
Sturm}

\begin{problem}\label{pb:b2:diffeq:1}
Satu persamaan menguasai dunia mekanika dan kelistrikan:
\[
x'' + 2\zeta\omega\,x' + \omega^2 x = F(t),
\qquad \omega > 0,\ \zeta \geq 0 .
\]
Soal ini mengkajinya selengkapnya --- lewat penggolongan
trace--determinan bagi sistem linear bidang,
ketiga rezim peredamannya, tanggapan keadaan tunak terhadap
pemaksaan berkala beserta puncak \emph{resonansi}\index{resonansi}nya
dan malapetaka \omterm{pb:b2:diffeq:1}{resonansinya} --- lalu meninggalkan koefisien
tetapnya demi \emph{teorema pemisahan dan pembandingan
Sturm}\index{teorema pembandingan Sturm}, yang mengendalikan
nol penyelesaian $y'' + q(t)y = 0$ tanpa satu rumus pun.

\medskip
\noindent\textbf{Bagian I --- Bidang trace--determinan.} Misalkan
$A \in \mathcal M_2(\R)$, $\tau = \operatorname{tr}A$, $\delta
= \det A$, dan $\Delta = \tau^2 - 4\delta$.
\begin{enumerate}
  \item Tunjukkan bahwa \omterm{def:b2:reduction:eigen}{nilai eigen} $A$ adalah
        $\frac{\tau\pm\sqrt\Delta}{2}$ lalu golongkanlah: dua \omterm{def:b2:reduction:eigen}{nilai
        eigen} real bertanda berlawanan bila dan hanya bila $\delta < 0$; dua \omterm{def:b2:reduction:eigen}{nilai
        eigen} real bertanda sama bila dan hanya bila $\delta > 0$ dan $\Delta
        \geq 0$ (dengan tanda $\tau$); dan sepasang sekawan tak real bila dan hanya bila
        $\Delta < 0$ (dengan bagian real $\frac\tau2$).
  \item (Pelana, $\delta < 0$) Dengan \omterm{def:b2:reduction:eigen}{nilai eigen} $\mu < 0 <
        \lambda$ dan \omterm{def:b2:reduction:eigen}{vektor eigen} $v_\pm$, tulislah penyelesaian
        umumnya lalu paparkan lintasannya: yakni dua sinar stabil dan
        dua sinar yang takstabil, sedangkan orbit lainnya asimtotik terhadap
        keduanya. Mengapa tak ada penyelesaian selain $0$ yang dapat tetap terbatas pada seluruh
        $\R$?
  \item (Simpul, $\delta > 0$, $\Delta > 0$) Untuk $\mu < \lambda
        < 0$: tunjukkan bahwa setiap penyelesaian tak nol menuju $0$ dan
        bahwa semua orbitnya kecuali yang berada di sumbu cepatnya tiba
        \emph{secara menyinggung arah eigen yang lambat}
        \emph{(bandingkan $\eu^{\mu t}$ dan $\eu^{\lambda t}$)}.
  \item (Spiral dan pusat, $\Delta < 0$) Dengan menulis
        \omterm{def:b2:reduction:eigen}{nilai eigennya} $\alpha \pm \iu\beta$, pakailah
        \cref{exo:b2:diffeq:5} (setelah penggantian basis yang real,
        yang diterima tanpa bukti dalam keumuman itu atau dibuktikan bagi sistem
        Bagian II, yakni yang dipakai di bawah) untuk memaparkan
        orbitnya: yakni spiral yang konvergen bagi $\alpha =
        \frac\tau2 < 0$, yang divergen bagi $\tau > 0$, dan kurva
        tertutup (yakni pusat) bagi $\tau = 0$.
  \item (Kasus perbatasannya) Untuk \omterm{def:b2:reduction:eigen}{nilai eigen} ganda (yakni $\Delta =
        0$): tunjukkan $\eu^{tA} = \eu^{\lambda t}(I + tN)$ dengan $N
        = A - \lambda I$ yang nilpoten, lalu bedakanlah bintangnya
        ($N = 0$) dari simpul takwajarnya ($N \neq 0$).
        Rangkumlah Bagian I dalam gambaran trace--determinan pada
        gambar bab ini.
\end{enumerate}

\medskip
\noindent\textbf{Bagian II --- Osilator teredam.} Kini $F =
0$: $x'' + 2\zeta\omega x' + \omega^2x = 0$, yakni $X' = AX$
dengan $A = \begin{pmatrix} 0 & 1\\ -\omega^2 &
-2\zeta\omega\end{pmatrix}$.
\begin{enumerate}[resume]
  \item Hitunglah $\tau, \delta, \Delta$ lalu tempatkanlah ketiga
        rezimnya pada bidang trace--determinannya:
        \emph{teredam kurang} $0 < \zeta < 1$ (yakni spiral stabil),
        \emph{teredam kritis} $\zeta = 1$ (yakni \omterm{def:b2:reduction:eigen}{nilai eigen}
        ganda), \emph{teredam lebih} $\zeta > 1$ (yakni simpul
        stabil); sedangkan $\zeta = 0$ adalah pusatnya.
  \item Selesaikan ketiga rezimnya secara gamblang:
        \[
        \zeta < 1:\ \eu^{-\zeta\omega t}\bigl(a\cos\omega_d t
        + b\sin\omega_dt\bigr),\ \omega_d =
        \omega\sqrt{1-\zeta^2};
        \qquad
        \zeta = 1:\ (a + bt)\,\eu^{-\omega t};
        \]
        untuk $\zeta > 1$: dua eksponensial real. Definisikanlah
        pseudokalanya $\frac{2\pi}{\omega_d}$ lalu tunjukkan bahwa
        nisbah maksimum berurutan $\abs x$ adalah konstanta
        $\eu^{-2\pi\zeta/\sqrt{1-\zeta^2}}$ (yakni dekremen
        logaritmiknya).
  \item (Asas penutup pintu) Untuk $\zeta \geq 1$, laju
        peluruhannya dikendalikan \omterm{def:b2:reduction:eigen}{nilai eigen} yang paling lambat
        $\lambda_{\mathrm{lambat}} = -\omega\bigl(\zeta -
        \sqrt{\zeta^2-1}\bigr)$. Tunjukkan bahwa
        $\abs{\lambda_{\mathrm{lambat}}} =
        \frac{\omega}{\zeta + \sqrt{\zeta^2 - 1}}$ merupakan
        fungsi \emph{turun} atas $\zeta \geq 1$: jadi
        peredaman kritis $\zeta = 1$ memberi kepulangan ke
        istirahat yang tercepat tanpa berayun.
  \item (Energi) Misalkan $E(t) = \frac12x'^2 +
        \frac12\omega^2x^2$. Buktikan $E' = -2\zeta\omega\,x'^2
        \leq 0$, lalu turunkan bahwa untuk $\zeta > 0$ persamaannya
        tak punya penyelesaian berkala yang tak nol \emph{(sebab sebuah kala akan
        memaksa $E$ tetap, sehingga $x' \equiv 0$)}.
  \item Jelaskan dalam dua kalimat mengapa pusat $\zeta = 0$
        bersifat \emph{rapuh secara struktural}: sebab sembarang $\zeta > 0$,
        sekecil apa pun, menghancurkan kekalaannya --- dan di mana
        hal itu tampak pada bidang trace--determinannya (yakni
        garis pusatnya berinterior kosong).
\end{enumerate}

\medskip
\noindent\textbf{Bagian III --- Osilasi terpaksa dan
\omterm{pb:b2:diffeq:1}{resonansi}.} Kini $F(t) = F\cos(\gamma t)$ dengan $F, \gamma > 0$.
\begin{enumerate}[resume]
  \item ($\zeta > 0$: keadaan tunaknya) Carilah $x_p =
        \Re\bigl(z\,\eu^{\iu\gamma t}\bigr)$: lalu tunjukkan
        \[
        z = \frac{F}{\omega^2 - \gamma^2 +
        2\iu\zeta\omega\gamma},
        \qquad
        A(\gamma) := \abs z = \frac{F}{\sqrt{(\omega^2 -
        \gamma^2)^2 + 4\zeta^2\omega^2\gamma^2}} ,
        \]
        lalu tulislah $x_p = A(\gamma)\cos(\gamma t - \varphi)$
        dengan $\tan\varphi =
        \frac{2\zeta\omega\gamma}{\omega^2-\gamma^2}$.
  \item Tunjukkan bahwa \emph{setiap} penyelesaiannya adalah $x_p$ ditambah
        sebuah transien dari Bagian II, yang menuju $0$: jadi apa pun
        data awalnya, sistemnya mengunci pada keadaan
        tunaknya --- dengan amplitudo $A(\gamma)$ dan ketertinggalan fasa
        $\varphi$.
  \item (Kurva \omterm{pb:b2:diffeq:1}{resonansinya}) Maksimumkan $A$: tunjukkan bahwa
        $A(\gamma)$ punya maksimum dalam bila dan hanya bila $\zeta <
        \frac{1}{\sqrt2}$, yakni di
        \[
        \gamma_* = \omega\sqrt{1 - 2\zeta^2},
        \qquad
        A(\gamma_*) =
        \frac{F}{2\zeta\omega^2\sqrt{1-\zeta^2}} ,
        \]
        dan bahwa untuk $\zeta$ yang kecil puncaknya memperkuat
        tanggapan statisnya $A(0) = \frac F{\omega^2}$ sebesar
        faktor $\approx \frac{1}{2\zeta}$.
  \item ($\zeta = 0$, di luar \omterm{pb:b2:diffeq:1}{resonansi}) Untuk $\gamma \neq \omega$,
        tunjukkan bahwa penyelesaian dengan $x(0) = x'(0) = 0$ adalah
        \[
        x(t) = \frac{F}{\omega^2 -
        \gamma^2}\bigl(\cos\gamma t - \cos\omega t\bigr)
        = \frac{2F}{\omega^2-\gamma^2}
        \sin\frac{(\omega-\gamma)t}{2}
        \sin\frac{(\omega+\gamma)t}{2} :
        \]
        yang terbatas, dengan \emph{layangan} --- yakni ayunan cepat
        di bawah selubung yang lambat --- ketika $\gamma$ dekat ke
        $\omega$.
  \item ($\zeta = 0$, pada \omterm{pb:b2:diffeq:1}{resonansi}) Untuk $\gamma = \omega$, tunjukkan
        bahwa $x_p(t) = \frac{F}{2\omega}\,t\sin(\omega t)$ merupakan
        sebuah penyelesaian, lalu perolehlah ia sebagai limit pertanyaan
        14 ketika $\gamma \to \omega$: jadi amplitudonya tumbuh
        linear selamanya --- itulah malapetaka \omterm{pb:b2:diffeq:1}{resonansinya}.
  \item (Kaitan Fourier) Sebuah pemaksaan berkala yang umum terurai
        menjadi harmonik (lewat bab Fourier); jadi berkat kelinearannya
        keadaan tunaknya adalah jumlah tanggapan harmoniknya. Untuk
        sebuah osilator tak teredam berfrekuensi $\omega = 3$ yang
        dipaksa gelombang persegi bertipe
        \cref{exo:b2:fourier:1} (yakni berharmonik pada setiap
        bilangan bulat ganjil), harmonik mana yang beresonansi? Satu kalimat tentang
        mengapa para insinyur takut pada gelombang persegi.
\end{enumerate}

\medskip
\noindent\textbf{Bagian IV --- Teorema Sturm.} Tinjaulah $y''
+ q(t)\,y = 0$ pada sebuah interval $I$, dengan $q$ \omterm{def:b2:metric:continuity}{kontinu}. (Setiap
persamaan $y'' + ay' + by = 0$ menyusut menjadi bentuk normal ini lewat
penyulihan $y = u\exp\bigl(-\frac12\int a\bigr)$; dan
pertanyaan 21 menunjukkan sebuah varian siasat itu beraksi.)
\begin{enumerate}[resume]
  \item Untuk dua penyelesaian $y_1, y_2$, tunjukkan bahwa \omterm{def:b2:diffeq:wronskian}{wronskiannya}
        $W = y_1y_2' - y_1'y_2$ bersifat \emph{tetap}, yang nol bila dan hanya bila
        penyelesaiannya sebanding; dan bahwa penyelesaian tak nol
        hanya punya nol yang \emph{sederhana dan terpencil}.
  \item (Pemisahan Sturm) Misalkan $y_1, y_2$ penyelesaian yang
        bebas dan $a < b$ dua nol berurutan milik $y_1$.
        Buktikan bahwa $y_2$ lenyap tepat sekali di
        $\intoo{a}{b}$ \emph{(nilailah konstanta $W$ di $a$
        dan $b$: sebab $W = y_1'y_2$ di sana, dan $y_1'(a)$ serta
        $y_1'(b)$ bertanda berlawanan)}: jadi nol penyelesaian yang
        bebas saling menyisip.
  \item (Pembandingan Sturm) Misalkan $q_1 \leq q_2$ pada $I$, dengan $y \neq
        0$ memenuhi $y'' + q_1y = 0$, $z \neq 0$ memenuhi $z'' + q_2z
        = 0$, dan $a < b$ nol berurutan milik $y$. Tunjukkan bahwa
        $z$ lenyap di $\intcc{a}{b}$ --- dan tegas di dalamnya bila
        $q_1 < q_2$ di suatu tempat pada $\intoo ab$ \emph{(bila $z \neq
        0$ pada $\intoo ab$, kajilah $(yz' - y'z)' = (q_1 -
        q_2)yz$ dengan tanda $y, z$ yang tetap lalu bandingkan
        nilai perbatasannya)}.
  \item Turunkan \emph{batas jaraknya}: bahwa bila $0 < m^2 \leq
        q(t) \leq M^2$ pada $I$, maka sembarang dua nol berurutan
        $a < b$ milik penyelesaian tak nol $y'' + qy = 0$
        memenuhi
        \[
        \frac{\pi}{M} \;\leq\; b - a \;\leq\; \frac{\pi}{m}
        \]
        \emph{(bandingkan dengan $u'' + M^2u = 0$ dan $u'' + m^2u
        = 0$, yang nolnya berjarak $\frac\pi M$ dan
        $\frac\pi m$)}. Periksalah pada osilator harmoniknya.
  \item Transformasikan $ty'' + 2y' + ty = 0$
        (\cref{exo:b2:diffeq:7}) lewat $u = ty$ menjadi $u'' + u =
        0$, perolehlah kembali penyelesaiannya $\frac{\sin t}t$ dan
        $\frac{\cos t}{t}$ seketika, lalu simpulkan bahwa
        nol setiap penyelesaian tak nolnya berjarak tepat
        $\pi$: yakni pandangan dunia Sturm --- bahwa nolnya dikendalikan
        koefisien $q$, ada rumus atau tidak.
\end{enumerate}

\medskip
\noindent\textbf{Bagian V --- Duhamel dan perbatasan
keterbatasan.}
\begin{enumerate}[resume]
  \item (Duhamel bagi osilatornya) Tunjukkan bahwa untuk
        $F$ yang \omterm{def:b2:metric:continuity}{kontinu}, penyelesaian $x'' + \omega^2x =
        F(t)$ dengan $x(0) = x'(0) = 0$ adalah
        \[
        x(t) = \frac1\omega\int_0^t\sin\bigl(\omega(t -
        s)\bigr)F(s)\,\dd s ,
        \]
        lalu turunkan ulang penyelesaian resonan pertanyaan 15
        darinya dengan $F(s) = F\cos(\omega s)$
        \emph{(lewat hasil-kali-ke-jumlah)}.
  \item ($\zeta > 0$: masukan terbatas, keluaran terbatas) Tunjukkan
        bahwa untuk $\zeta > 0$ dan \emph{sembarang} $F$ \omterm{def:b2:metric:continuity}{kontinu}
        yang terbatas, setiap penyelesaian persamaan teredamnya bersifat terbatas
        pada $\intco{0}{\infty}$ \emph{(lewat variasi konstanta
        ditambah peluruhan eksponensial $\vertiii{\eu^{tA}} \leq
        C\eu^{-\alpha t}$ pada \cref{exo:b2:diffeq:8})}.
  \item ($\zeta = 0$) Tunjukkan bahwa tanpa peredaman, pemaksaan
        berkala yang terbatas menjaga semua penyelesaiannya tetap terbatas
        \emph{kecuali} persis pada \omterm{pb:b2:diffeq:1}{resonansinya} (yakni $\gamma =
        \omega$, pertanyaan 15 lawan pertanyaan 14): jadi peredamanlah
        yang mengubah perbatasan keterbatasannya menjadi kestabilan \omterm{def:b2:funcseq:def}{seragam}.
  \item Rangkuman. Satu kalimat untuk masing-masing: (i) bagaimana
        bidang trace--determinan menata Bagian I--II dan
        di mana pemaksaan (Bagian III) meninggalkannya; (ii) makna
        fisis $\gamma_*$, $A(\gamma_*)$ dan
        faktor $\frac1{2\zeta}$; (iii) apa yang dikatakan teorema
        Sturm yang tak dapat dikatakan rumus gamblang; (iv)
        dua hasil mana pada soal ini yang kelak dipakai ulang secara diam-diam
        oleh sisa buku ini (yakni \omterm{def:b2:diffeq:wronskian}{wronskian} bertetapan Liouville;
        dan kestabilan masukan-terbatas).

\end{enumerate}
\end{problem}
